% © 2026 Ing. David Jaroš — CC BY-NC-ND 4.0
%
% This work is licensed under a Creative Commons Attribution-NonCommercial-NoDerivatives
% 4.0 International License (CC BY-NC-ND 4.0).
%
% License History: Earlier drafts (up to v0.3) were released under CC BY 4.0.
% From v0.4 onward, all material is released under CC BY-NC-ND 4.0 to protect
% the integrity of the theoretical work during ongoing academic development.

\documentclass[12pt]{article}
\usepackage{amsmath,amssymb,amsthm}
\usepackage{geometry}
\usepackage{hyperref}
\usepackage{tcolorbox}
\usepackage{xcolor}
\geometry{margin=1in}

% Theorem-like environments
\newtheorem{definition}{Definition}
\newtheorem{proposition}{Proposition}
\newtheorem{remark}{Remark}
\newtheorem{conjecture}{Conjecture}

\title{ST-1: Imaginary Sector Objects in UBT\\
Formal Definition and Field-Theoretic Status\\[0.5em]
\large\textsc{Speculative Extension — Invisibility Track}}
\author{David Jaroš}
\date{2026-03-03}

\begin{document}
\maketitle

\begin{tcolorbox}[colback=yellow!10, title=Status]
\textbf{SPECULATIVE EXTENSION} — This content is not part of
the canonical UBT framework.  It is an exploratory investigation
into the possibility of field configurations that reside entirely
in the imaginary sector of the biquaternionic metric, and has not
been validated against the canonical field equations.
\end{tcolorbox}

\vspace{1em}
\noindent\textit{Output of Sub-task ST-1 from the Invisibility / Imaginary Sector
investigation track.  See also:
\texttt{st2\_sector\_interaction.tex},
\texttt{st3\_stability.tex},
\texttt{st4\_ctc\_invisibility\_connection.tex},
\texttt{st5\_falsifiability.md},
and the CTC files in \texttt{speculative\_extensions/causality/}.}

%----------------------------------------------------------------------
\section*{Abstract}
%----------------------------------------------------------------------
The UBT biquaternionic metric $\mathcal{G}_{\mu\nu} = g_{\mu\nu}
+ i\,h_{\mu\nu} + j\,a_{\mu\nu} + k\,b_{\mu\nu}$ admits
a natural decomposition into a \emph{real sector} ($g_{\mu\nu}$)
and an \emph{imaginary sector} ($h_{\mu\nu}$, $a_{\mu\nu}$,
$b_{\mu\nu}$).  We formally define an \emph{imaginary sector object}
as a field configuration $\Theta_\mathrm{imag}$ for which
$\operatorname{Re}(\mathcal{G}_{\mu\nu}[\Theta_\mathrm{imag}]) = 0$
while $\operatorname{Im}(\mathcal{G}_{\mu\nu}[\Theta_\mathrm{imag}]) \neq 0$.
Three foundational questions are addressed: (i) consistency of
the configuration $\Theta_\mathrm{imag} = i\,\Theta_\mathrm{real}$
with AXIOM A, (ii) well-definedness of the UBT action on this
configuration, and (iii) satisfaction of the UBT field equations.
All conclusions are \textbf{speculative}.

%----------------------------------------------------------------------
\section{Background: UBT Field and Metric Structure}
%----------------------------------------------------------------------

The canonical UBT $\Theta$-field takes values in the biquaternion
algebra $\mathcal{B} = \mathbb{C} \otimes \mathbb{H}$:
\begin{equation}
  \Theta(q,\tau) \;=\;
  \Theta_0 + \Theta_1 \mathbf{i} + \Theta_2 \mathbf{j} + \Theta_3 \mathbf{k},
  \qquad \Theta_n \in \mathbb{C},
  \label{eq:theta_field}
\end{equation}
where $\mathbf{i}, \mathbf{j}, \mathbf{k}$ are the quaternion
imaginary units and $q \in \mathbb{H}$, $\tau = t + i\psi \in
\mathbb{C}$ denote the biquaternionic spacetime coordinate and
complex time respectively.

The induced biquaternionic metric is
\begin{equation}
  \mathcal{G}_{\mu\nu}
  \;=\; g_{\mu\nu} + i\,h_{\mu\nu} + j\,a_{\mu\nu} + k\,b_{\mu\nu},
  \qquad g_{\mu\nu} = \operatorname{Re}(\mathcal{G}_{\mu\nu}),
  \label{eq:biquat_metric}
\end{equation}
and the UBT field equation is
\begin{equation}
  \mathcal{E}_{\mu\nu}[\Theta,\mathcal{G}]
  \;=\; \kappa\,\mathcal{T}_{\mu\nu}[\Theta,\mathcal{G}].
  \label{eq:ubt_field_eq}
\end{equation}
Standard GR is recovered by taking the real projection of
\eqref{eq:ubt_field_eq}, which yields the Einstein equations with
an effective stress-energy $T_{\mu\nu}^{\mathrm{eff}} =
\operatorname{Re}(\mathcal{T}_{\mu\nu})$.

%----------------------------------------------------------------------
\section{Definitions}
%----------------------------------------------------------------------

\begin{definition}[Real-sector object]
A field configuration $\Theta_\mathrm{real}$ is a
\emph{real-sector object} if the induced biquaternionic metric
satisfies
\begin{equation}
  \operatorname{Re}\!\left(\mathcal{G}_{\mu\nu}[\Theta_\mathrm{real}]\right)
  \;\neq\; 0
  \quad \text{on a non-empty open set of spacetime}.
\end{equation}
Such an object contributes to the physical metric $g_{\mu\nu}$
and is, in principle, directly detectable by real-sector observers.
\end{definition}

\begin{definition}[Imaginary-sector object]
\label{def:imag_obj}
A field configuration $\Theta_\mathrm{imag}$ is an
\emph{imaginary-sector object} if
\begin{align}
  \operatorname{Re}\!\left(\mathcal{G}_{\mu\nu}[\Theta_\mathrm{imag}]\right)
  &\;=\; 0
  \quad \text{everywhere (or asymptotically flat to zero),}
  \label{eq:imag_cond1}\\
  \operatorname{Im}\!\left(\mathcal{G}_{\mu\nu}[\Theta_\mathrm{imag}]\right)
  &\;\neq\; 0
  \quad \text{on a non-empty open set.}
  \label{eq:imag_cond2}
\end{align}
Such an object contributes \emph{only} to the imaginary metric
components $\{h_{\mu\nu}, a_{\mu\nu}, b_{\mu\nu}\}$ and is
\emph{a priori} invisible to observers whose apparatus couples
solely to $g_{\mu\nu}$.
\end{definition}

\begin{remark}
Definition~\ref{def:imag_obj} is a \emph{formal} definition.
Whether a physically consistent realisation exists is the subject
of the questions below.  The word ``invisible'' here means
invisible to the real-sector metric; imaginary-sector objects
may still be detectable through the back-reaction mechanism
discussed in \texttt{st2\_sector\_interaction.tex}.
\end{remark}

%----------------------------------------------------------------------
\section{Question 1: Is $\Theta_\mathrm{imag} = i\,\Theta_\mathrm{real}$
Valid Under AXIOM A?}
%----------------------------------------------------------------------

AXIOM A of UBT (following the notation of the canonical core document)
states that the $\Theta$-field is a smooth section of the biquaternion
bundle $\mathcal{B}$ over the spacetime manifold, subject to the
condition that the induced real metric $g_{\mu\nu}$ has Lorentzian
signature $(-,+,+,+)$.

\paragraph{The $i$-rotation map.}
Consider the map $\Theta \mapsto i\,\Theta$ where $i$ is the
complex unit in $\mathcal{B} = \mathbb{C}\otimes\mathbb{H}$.
This maps $\Theta_\mathrm{real} \mapsto i\,\Theta_\mathrm{real}$.

Since $\mathcal{B}$ is a unital $\mathbb{C}$-algebra, $i\,\Theta$
is a valid element of $\mathcal{B}$ for any $\Theta \in \mathcal{B}$.
The map is smooth and bijective on $\mathcal{B}$.

\paragraph{Effect on the metric.}
The biquaternionic metric is bilinear in $\Theta$ (schematically):
\begin{equation}
  \mathcal{G}_{\mu\nu}[\Theta]
  \;\sim\; \partial_\mu\Theta^\dagger \partial_\nu\Theta + \cdots
\end{equation}
Under $\Theta \to i\,\Theta$:
\begin{equation}
  \mathcal{G}_{\mu\nu}[i\,\Theta]
  \;\sim\; (i\,\partial_\mu\Theta)^\dagger (i\,\partial_\nu\Theta)
  \;=\; (i^*)\,(\partial_\mu\Theta^\dagger)\,(i)\,(\partial_\nu\Theta)
  \;=\; (-i)(i)\,\partial_\mu\Theta^\dagger \partial_\nu\Theta
  \;=\; \mathcal{G}_{\mu\nu}[\Theta].
\end{equation}
Here we used the Hermitian conjugation rule $(i\,A)^\dagger = A^\dagger(-i)$
for $i \in \mathbb{C}$.  Consequently, if the metric depends on $\Theta$
only through the combination $\partial\Theta^\dagger \partial\Theta$, then
$\mathcal{G}_{\mu\nu}[i\,\Theta] = \mathcal{G}_{\mu\nu}[\Theta]$ and the
imaginary rotation does \emph{not} produce a new configuration — it leaves
the metric unchanged.

\paragraph{Non-bilinear metric contributions.}
If the metric contains terms linear in $\Theta$ (or in $\partial\Theta$
without the Hermitian conjugate), the $i$-rotation maps real parts to
imaginary parts:
\begin{equation}
  \mathcal{G}_{\mu\nu}[i\,\Theta]
  \;\ni\; i\,(\text{linear terms in }\Theta)
  \;\Rightarrow\;
  \operatorname{Re}(\mathcal{G}_{\mu\nu}[i\,\Theta])
  \;\ni\; -\operatorname{Im}(\text{linear terms in }\Theta).
\end{equation}
In this case, starting from a $\Theta_\mathrm{real}$ with
$\operatorname{Im}(\partial\Theta_\mathrm{real}) = 0$, the rotated
field $i\,\Theta_\mathrm{real}$ could have
$\operatorname{Re}(\mathcal{G}_{\mu\nu}) = 0$.

\begin{conjecture}[AXIOM A consistency]
\label{conj:axiom_a}
The configuration $\Theta_\mathrm{imag} = i\,\Theta_\mathrm{real}$ is
consistent with AXIOM A (smoothness and biquaternion membership) but
may violate the subsidiary condition that $g_{\mu\nu}$ has Lorentzian
signature, because condition \eqref{eq:imag_cond1} forces $g_{\mu\nu} = 0$,
which is degenerate.  Whether the axiom system permits degenerate metrics
in an intermediate formal sense (before projection to physical observables)
is an open question.
\end{conjecture}

%----------------------------------------------------------------------
\section{Question 2: Is the UBT Action Well-Defined on
$\Theta_\mathrm{imag}$?}
%----------------------------------------------------------------------

The UBT action takes the form (schematically)
\begin{equation}
  S[\Theta, \mathcal{G}]
  \;=\; \int_{\mathcal{M}} \mathcal{L}(\Theta, \partial\Theta, \mathcal{G})
  \,\sqrt{-\det(\operatorname{Re}\,\mathcal{G})}\;d^4x\,d\psi,
  \label{eq:action}
\end{equation}
where the volume element $\sqrt{-\det g}$ involves the \emph{real}
metric determinant.

\paragraph{Degeneracy of the volume element.}
For an imaginary-sector object satisfying \eqref{eq:imag_cond1},
$g_{\mu\nu} = 0$ everywhere, so $\det g = 0$ and the volume element
$\sqrt{-\det g} = 0$.  The action \eqref{eq:action} is then
identically zero.

\paragraph{Alternative: action over complex volume.}
One could formally define a complexified action
\begin{equation}
  S_\mathbb{C}[\Theta, \mathcal{G}]
  \;=\; \int_{\mathcal{M}} \mathcal{L}(\Theta, \partial\Theta, \mathcal{G})
  \,\sqrt{-\det(\mathcal{G})}\;d^4x\,d\psi,
\end{equation}
where $\det(\mathcal{G})$ is the determinant of the full biquaternionic
metric.  This is mathematically well-defined when
$\operatorname{Im}(\mathcal{G}_{\mu\nu}) \neq 0$, even if
$\operatorname{Re}(\mathcal{G}_{\mu\nu}) = 0$.

\begin{remark}
Using $S_\mathbb{C}$ requires extending the variational principle to
complex field variations, which is non-standard.  The resulting
field equations may not be real, raising questions about the physical
interpretation of the imaginary parts of the equations of motion.
This extension is not part of the canonical UBT framework.
\end{remark}

%----------------------------------------------------------------------
\section{Question 3: Do Imaginary-Sector Objects Satisfy the
Field Equations?}
%----------------------------------------------------------------------

The real projection of the UBT field equation \eqref{eq:ubt_field_eq}
yields
\begin{equation}
  R_{\mu\nu} - \tfrac{1}{2}g_{\mu\nu}R
  \;=\; 8\pi G\,T_{\mu\nu}^{\mathrm{eff}},
\end{equation}
where $T_{\mu\nu}^{\mathrm{eff}} = \operatorname{Re}(\mathcal{T}_{\mu\nu})$.
The imaginary projection yields a separate set of equations for the
imaginary metric components.

For $\Theta_\mathrm{imag}$ with $g_{\mu\nu} = 0$:
\begin{itemize}
  \item The Einstein tensor $G_{\mu\nu} = R_{\mu\nu} - \frac{1}{2}g_{\mu\nu}R$
    is formally zero (as $g_{\mu\nu}$ and all its derivatives vanish), but
    the Christoffel symbols $\Gamma^\lambda_{\mu\nu} \propto g^{\lambda\sigma}
    \partial g_{\sigma\mu}/\partial x^\nu$ are ill-defined when $g_{\mu\nu} = 0$
    because the inverse $g^{\mu\nu}$ diverges.
  \item The matter field equation $\nabla^2\Theta - \partial V/\partial\Theta^\dagger = 0$
    involves the covariant d'Alembertian $\nabla^2 = g^{\mu\nu}\nabla_\mu\nabla_\nu$,
    which is also ill-defined when $g_{\mu\nu} = 0$.
\end{itemize}

\begin{proposition}[Field equation obstruction]
\label{prop:obstruction}
A globally imaginary-sector object satisfying $g_{\mu\nu} \equiv 0$
everywhere cannot satisfy the canonical UBT field equations
\eqref{eq:ubt_field_eq} in the real-projection form, because the
real metric sector is degenerate and the associated differential
operators are ill-defined.
\end{proposition}

\paragraph{Localised near-imaginary objects.}
Proposition~\ref{prop:obstruction} does not exclude configurations
that are ``almost imaginary'' — i.e., have $g_{\mu\nu} = \epsilon\,
\bar{g}_{\mu\nu}$ with $\epsilon \ll 1$, so that the metric is
nearly degenerate but the field equations remain well-defined.  In
the limit $\epsilon \to 0$ one approaches the imaginary sector.

Such \emph{near-imaginary} configurations may provide a physically
meaningful approximation to the ideal imaginary-sector object,
and their interaction with the real sector is discussed in
\texttt{st2\_sector\_interaction.tex}.

%----------------------------------------------------------------------
\section{Summary}
%----------------------------------------------------------------------
\begin{enumerate}
  \item The $i$-rotation $\Theta \to i\,\Theta$ is algebraically
    valid in $\mathcal{B} = \mathbb{C}\otimes\mathbb{H}$, but for
    a metric that is Hermitian-bilinear in $\Theta$ the rotation
    leaves $\mathcal{G}_{\mu\nu}$ unchanged.  A non-trivial
    imaginary-sector object requires a metric with linear or
    non-Hermitian dependence on $\Theta$.
  \item The standard UBT action is ill-defined on a configuration
    with $g_{\mu\nu} = 0$ because the volume element vanishes.
    A complexified action is formally well-defined but non-standard
    and outside the canonical UBT framework.
  \item The canonical field equations cannot be satisfied globally
    on an imaginary-sector configuration because the covariant
    differential operators require a non-degenerate $g_{\mu\nu}$.
  \item The physically relevant generalisation is a
    \emph{near-imaginary} configuration with $|g_{\mu\nu}| \ll
    |h_{\mu\nu}|$, which is tractable in perturbation theory and
    examined in \texttt{st2\_sector\_interaction.tex}.
\end{enumerate}

\medskip
\noindent\textbf{All statements in this document are speculative.
No claim constitutes an established result of the canonical UBT framework.}

%----------------------------------------------------------------------
\bibliographystyle{plain}
\begin{thebibliography}{9}
\bibitem{UBT_core}
  D.~Jaroš,
  \emph{Unified Biquaternion Theory: core field equations},
  Repository document \texttt{consolidation\_project/ubt\_core\_main.tex} (2025).
\bibitem{CTC_causal}
  D.~Jaroš,
  \emph{ST-1: Causal Structure of Complex Time in UBT},
  Repository document \texttt{speculative\_extensions/causality/complex\_time\_causal\_structure.tex}
  (2026).
\bibitem{CTC_conditions}
  D.~Jaroš,
  \emph{ST-3: Closed Timelike Curve Conditions in the Biquaternionic Metric},
  Repository document \texttt{speculative\_extensions/causality/ctc\_conditions.tex}
  (2026).
\end{thebibliography}

\end{document}
